\section{Variational state preparation}
\label{sec:ansatz}

\subsection{Illustrative start and end points}
\label{sec:points}

For each phase we define an illustrative start/end pair on the
$\tilde{Z}_{\mathcal R}$ map (Fig.~\ref{fig:D05}). The end points are
classical reference targets chosen deep inside the corresponding phase;
they are demonstration endpoints for the variational protocol, not the
actual final evolution endpoints. Keeping the reference points away from
the phase boundaries reduces sensitivity to finite-size boundary drift.
Start points share the dot style in three colors; end markers differ in
both shape and color, with arrows showing the reference direction of
each pair.

\begin{figure}[htbp]
\centering
\includegraphics[width=\columnwidth]{exp01_D05.pdf}
\caption{Illustrative start/end pairs on the $\tilde{Z}_{\mathcal R}$
phase diagram of Fig.~\ref{fig:D01}, one pair per phase with the
reference end point deep inside it; colors distinguish pairs and arrows
mark reference direction.
\label{fig:D05}}
\end{figure}

\subsection{Variational circuit structure}

Our variational circuits follow the variational eigenvalue solver
paradigm~\cite{peruzzo2014variational}: prepare an initial state, then
apply a parameterized ansatz. The classical loop minimizes
\begin{equation}
E(\theta)=\langle\psi(\theta)|H|\psi(\theta)\rangle,
\end{equation}
where $|\psi(\theta)\rangle$ is the circuit output at parameters
$\theta$. Variational algorithms of this form are surveyed in
Refs.~\cite{cerezo2021variational,tilly2022variational}, which collect
cost-function design, ansatz families, and best practices.
In practice the optimization landscape, not the
expressivity, is the bottleneck: gradients vanish exponentially with
qubit number in random circuits forming
$2$-designs~\cite{mcclean2018barren}, hardware-efficient circuits exhibit
abrupt trainability transitions with depth~\cite{campos2021abrupt}, while
optimal parameters can concentrate inverse-polynomially in problem
size~\cite{akshay2021parameter}, smooth Hamiltonian-variational solutions
transfer across sizes with non-vanishing
gradients~\cite{mele2022avoiding}, and barren-plateau regimes can hide
exponentially many traps~\cite{nemkov2025barren}. Krylov-like iterative
solvers offer a complementary route to ground
states~\cite{bharti2021iterative}. Hardware-efficient variational
optimization was first demonstrated on superconducting hardware for
small molecules and quantum
magnets~\cite{kandala2017hardware}, and spin-chain variational simulation
has since been pushed to $102$ qubits on cloud superconducting
devices~\cite{yu2023simulating}. Our orbit ansatz differs from all of
these in tying parameters by reflection symmetry and fixing the
entanglement structure at initialization, as described next.
The three initial states (odd-bond singlet
product and its counterparts, all in the $P=-2$ sector, where $P$ labels
the combined symmetry eigenvalue) pair with an
orbit ansatz in which mirror-symmetric bond pairs share rotation angles
(odd bonds $O_j\leftrightarrow O_{M+1-j}$, even bonds
$E_j\leftrightarrow E_{M-j}$). That is, the ansatz is built in two steps:
first fix the entanglement structure with a phase-matched initial state,
then dress all bonds with rotations tied by the reflection symmetry, so
the parameter count grows with the number of bond orbits rather than
bonds. One schematic per phase is shown in Fig.~\ref{fig:D06}(a)--(c): initial state
plus required ansatz circuit for the trivial, topological, and
AFM cases.
The construction keeps the trial states in the target symmetry sector.
The quality of the resulting states is evaluated later by the
hardware-simulator benchmarks.

\begin{figure}[htbp]
\centering
\includegraphics[width=\columnwidth]{exp04_D06a.pdf}
\caption{(a) Trivial initial state plus required ansatz circuit. The dashed
box $U^{(k)}$ is the $k$th variational layer; $\theta^{(k)}_j$
($\phi^{(k)}_j$) is the $XX{+}YY$ ($ZZ$) angle shared by the $j$th
mirror-symmetric bond orbit in layer $k$ (odd $O_j\leftrightarrow
O_{M+1-j}$, even $E_j\leftrightarrow E_{M-j}$; $\theta=\phi$ at
$\delta=0$). Even-bond sublayer first. \label{fig:D06}}
\end{figure}

\begin{figure}[htbp]
\centering
\addtocounter{figure}{-1}
\includegraphics[width=\columnwidth]{exp04_D06b.pdf}
\caption{(b) Topological initial state plus required ansatz circuit, same
notation as (a); odd-bond sublayer first.}
\end{figure}

\begin{figure}[htbp]
\centering
\addtocounter{figure}{-1}
\includegraphics[width=\columnwidth]{exp04_D06c.pdf}
\caption{(c) AFM initial state plus required ansatz circuit, same notation
as (a); even-bond sublayer first.}
\end{figure}
